\section*{Chapitre \ref{ch:pl:columnar-formats} --- Formats en colonnes}

\begin{solution}{exo:pl:columnar-formats:1}
Orienté lignes~: chaque ligne est lue en entier, $5 \times 8 \times 10^6 = 40$~Mo. En colonnes~: une seule colonne, $8 \times 10^6 = 8$~Mo.
\end{solution}

\begin{solution}{exo:pl:columnar-formats:2}
Bitmap~: 1\,000 bits, 125 octets. Décalages~: $1\,001 \times 4 = 4\,004$ octets. Données~: 49 octets. Au total 4\,178 octets~; les valeurs nulles coûtent
leur bit et un décalage chacun, sans donnée.
\end{solution}

\begin{solution}{exo:pl:columnar-formats:3}
Quatre valeurs demandent 2 bits chacune~: $50\,000 \times 2 / 8 = 12\,500$ octets d'indices. \texttt{firm.colfile} ajoute un en-tête de 9 octets
(nombre d'entrées, taille du dictionnaire, largeur en bits) et le dictionnaire lui-même stocké en clair, 40 octets (un compteur, cinq décalages de 4 octets et 16
caractères)~: 12\,549.
\end{solution}

\begin{solution}{exo:pl:columnar-formats:4}
Dans l'ordre de capture, les cotations consécutives des cinquante symboles sont espacées d'environ une demi-seconde, soit $4.7\times10^8$~ns, si bien que chaque delta demande
environ 29 bits, et la largeur de chaque bloc est fixée par son plus grand delta~: 32.2 bits par valeur mesurée. Triées par symbole, les lignes consécutives
sont des cotations d'un même symbole, espacées d'environ 23 secondes, et la série de chaque symbole commence par un saut en arrière dans le temps~: 44.9 bits par valeur. 
L'\omterm{def:pl:columnar-formats:encodings}{encodage delta} paie la taille des pas, et trier selon une autre clé rend les pas de temps plus grands.
\end{solution}

\begin{solution}{exo:pl:columnar-formats:5}
Chaque \omterm{def:pl:columnar-formats:rowgroup}{groupe de lignes} d'un fichier trié par temps contient des cotations des cinquante symboles, si bien que sa \omterm{def:pl:columnar-formats:pushdown}{carte de zones} sur le symbole va de S000 à S049 et qu'aucun groupe ne peut
être ignoré. Le lecteur applique aussi la projection~: il ne lit que la colonne des symboles (pour filtrer) et les cours acheteur et vendeur, soit 28~\%
du fichier, pour renvoyer 2~\% des lignes.
\end{solution}

\begin{solution}{exo:pl:columnar-formats:6}
À droite, les grands groupes font couvrir à chaque \omterm{def:pl:columnar-formats:pushdown}{carte de zones} une longue plage de temps, si bien qu'une part plus grande de chaque groupe lu se trouve hors de l'heure visée
(6.00~Mo pour 524\,288 lignes en ordre temporel). À gauche, le pied de fichier, que chaque requête lit avant tout le reste, croît avec le
nombre de groupes et de colonnes~: 444~ko pour 2\,048 lignes contre 115~ko pour 8\,192. Le minimum se trouve là où des \omterm{def:pl:columnar-formats:pushdown}{cartes de zones} grossières cessent de coûter
plus que les métadonnées.
\end{solution}

\begin{solution}{exo:pl:columnar-formats:7}
Le fichier compte 123 \omterm{def:pl:columnar-formats:rowgroup}{groupes de lignes}~; les \omterm{def:pl:columnar-formats:pushdown}{cartes de zones} sur le symbole et sur l'horodatage n'en retiennent que deux susceptibles de correspondre, et le lecteur ne décode que
les \omterm{def:pl:columnar-formats:rowgroup}{blocs de colonne} de ces deux groupes pour les quatre colonnes. Il renvoie 153 lignes, identique à la réponse de pyarrow sur les mêmes données.
\end{solution}

\begin{solution}{exo:pl:columnar-formats:8}
Les statistiques d'une colonne ne permettent d'élaguer que si ses valeurs sont resserrées dans chaque \omterm{def:pl:columnar-formats:rowgroup}{groupe de lignes}. Dans un fichier trié par temps, chaque groupe contient chaque
symbole, si bien que les \omterm{def:pl:columnar-formats:pushdown}{cartes de zones} du symbole couvrent tout et qu'une requête sur un seul symbole lit les colonnes projetées de chaque groupe (28~\% du
fichier du chapitre pour une requête renvoyant 2~\% des lignes). Trier par symbole à l'intérieur de chaque partition par date, ou conserver une seconde copie.
\end{solution}

\begin{solution}{pb:pl:columnar-formats:1}
\begin{enumerate}
\item Tout, puisque chaque ligne est lue en entier~; une \omterm{def:pl:columnar-formats:layouts}{disposition en colonnes} lirait $2/9 = 22\%$.
\item Un \omterm{def:pl:columnar-formats:validity}{bitmap de validité}, un tampon de décalages de $n+1$ entiers et un tampon de données.
\item C'est un bit du \omterm{def:pl:columnar-formats:validity}{bitmap de validité}~; la valeur elle-même n'occupe pas de place dans le tampon de données.
\item Dictionnaire (13 octets), dictionnaire (37\,917), delta (201\,085), dictionnaire (99\,097).
\item Trié, le symbole des 50\,000 lignes du jour forme 50 plages, stockées sous forme de 50 valeurs et de 50 longueurs.
\item Le minimum et le maximum d'une colonne dans le groupe~; cela permet au lecteur d'ignorer le groupe lorsqu'une condition du filtre
est fausse pour toute valeur de cette plage.
\item Non~: une ligne qui correspond se situe, pour chaque colonne filtrée, dans l'intervalle de son groupe, si bien que son groupe n'est jamais ignoré.
\item 86.0~\% pour 32\,768 lignes par groupe~; seule la projection fait gagner quelque chose, en écartant les colonnes que la requête ne nomme pas.
\item 3.30~\% pour un seul symbole et 3.13~\% pour tous les symboles.
\item 86.1~\%~: chaque groupe contient quelques symboles sur tout le mois, si bien que la plage temporelle de chaque groupe contient l'heure.
\item 6.00~Mo.
\item 0.39~Mo.
\item Le pied de fichier, lu en premier, énumère chaque \omterm{def:pl:columnar-formats:rowgroup}{bloc de colonne} des 489 groupes.
\item 3.9 fois (444 contre 115~ko).
\item 8\,192 sur ces données, d'après la mesure~; sur des données réelles, la même mesure avec le mélange de requêtes réel.
\item \textbf{Résultat nommé.} Avec des groupes de 8\,192 lignes, la requête sur un symbole et une heure lit 87.2~\% du fichier mélangé (17.0 sur
19.5~Mo)~, 2.65~\% en ordre temporel (0.39~Mo) et 2.79~\% trié par symbole puis par temps (0.36~Mo)~; des groupes plus petits comme plus grands coûtent
davantage (0.64~Mo pour 2\,048 lignes, 6.00~Mo pour 524\,288 en ordre temporel).
\item L'ordre temporel sert les requêtes d'une heure, l'ordre par symbole les requêtes d'un symbole~; un fichier mélangé ne sert ni l'un ni l'autre.
\item Partitionner par date, trier chaque jour par symbole puis par temps, choisir la taille des \omterm{def:pl:columnar-formats:rowgroup}{groupes de lignes} par la mesure~; conserver une copie triée par temps si
les coupes transversales comptent autant.
\item Le mélange réel des requêtes, la distribution réelle des symboles et des rythmes de cotation (quelques symboles dominent), et les octets lus, et non le
temps d'exécution seul.
\item À quel point les \omterm{def:pl:columnar-formats:pushdown}{cartes de zones} de ses colonnes de filtre sont resserrées dans l'ordre du fichier, et combien de colonnes il projette.
\end{enumerate}
\end{solution}

\begin{solution}{iq:pl:columnar-formats:1}
La recherche lit quelques champs de nombreux enregistrements, et les colonnes d'un même type se compressent bien et se vectorisent~; un fichier en colonnes ne lit que les
champs et les \omterm{def:pl:columnar-formats:rowgroup}{groupes de lignes} nécessaires. Un format en lignes convient mieux pour écrire les enregistrements au fil de leur arrivée et pour lire des enregistrements entiers
un à un~: la capture intrajournalière, l'état d'un ordre.
\iqlookfor{Lectures de quelques champs sur de nombreuses lignes~; ajouts et lectures ponctuelles favorisent les lignes.}
\end{solution}

\begin{solution}{iq:pl:columnar-formats:2}
Confronter le filtre au minimum et au maximum de chaque \omterm{def:pl:columnar-formats:rowgroup}{groupe de lignes} (ou de chaque \omterm{def:pl:columnar-formats:rowgroup}{page de données}) avant la lecture, et ignorer ceux qui ne peuvent
correspondre. Cela n'aide que si le fichier est trié (ou partitionné) pour que les colonnes filtrées aient des plages étroites dans chaque groupe.
\iqlookfor{\omterm{def:pl:columnar-formats:pushdown}{Cartes de zones}~; l'ordre de tri décide.}
\end{solution}

\begin{solution}{iq:pl:columnar-formats:3}
Horodatage~: \omterm{def:pl:columnar-formats:encodings}{encodage delta}, puisqu'il croît par petits pas. Symbole~: dictionnaire (peu de valeurs distinctes), ou \omterm{def:pl:columnar-formats:encodings}{codage par plages} une fois trié
par symbole. Prix~: dictionnaire si peu de prix distincts dans un \omterm{def:pl:columnar-formats:rowgroup}{bloc de colonne}, sinon \omterm{def:pl:columnar-formats:encodings}{encodage delta}, puisque les prix varient par pas de cotation.
\iqlookfor{Encodage adapté à la structure de la colonne, et dépendance à l'ordre de tri.}
\end{solution}

\begin{solution}{iq:pl:columnar-formats:4}
Un \omterm{def:pl:columnar-formats:validity}{bitmap de validité}, avec un bit par valeur, mis à 1 lorsque la valeur est présente~; une colonne sans valeur nulle peut s'en passer. Une valeur nulle coûte un bit (et, pour
les colonnes à largeur variable, un décalage), jamais une valeur sentinelle.
\iqlookfor{Bitmap et non valeur sentinelle~; ordre des bits, du poids faible au poids fort, comme détail.}
\end{solution}

\begin{solution}{iq:pl:columnar-formats:5}
Mesurer les octets lus, les \omterm{def:pl:columnar-formats:rowgroup}{groupes de lignes} lus et les colonnes lues par requête (un fichier de comptage, les statistiques du lecteur), à comparer aux lignes
renvoyées. Si les lectures dépassent largement le résultat, vérifier l'ordre de tri par rapport aux filtres, la taille des \omterm{def:pl:columnar-formats:rowgroup}{groupes de lignes}, et si
les statistiques sont écrites. Modifier alors la disposition (partitionner par date, trier selon le filtre principal, ajuster la taille des groupes) et mesurer de nouveau.
\iqlookfor{Octets lus rapportés aux lignes renvoyées, puis disposition, mesurée de nouveau.}
\end{solution}

\begin{solution}{iq:pl:columnar-formats:6}
Partitionner par date~; à l'intérieur de chaque jour, trier par symbole puis par temps, avec des \omterm{def:pl:columnar-formats:rowgroup}{groupes de lignes} dimensionnés par la mesure (quelques milliers de lignes) et
des statistiques sur le symbole et le temps~; encoder par dictionnaire les symboles et les conditions, encoder en delta les horodatages. Les historiques longs d'un seul symbole
lisent une seule petite plage par jour~; les coupes transversales d'une fenêtre temporelle lisent la partition du jour, projetée. Si les coupes transversales sont
fréquentes et volumineuses, ajouter une copie triée par temps.
\iqlookfor{\omterm{def:pl:capturing-and-storing-tick-data:partition}{Partitionnement}, plus ordre de tri, plus le constat qu'un seul ordre de tri ne peut servir les deux cas.}
\end{solution}
